\section{Related Work}
\label{sec:related}

\subsection{Grey wolf optimizer and its variants}

GWO was introduced by Mirjalili et al.\ \cite{mirjalili2014gwo} and models the
social hierarchy and encircling behaviour of grey wolves. Its simplicity and its
small number of parameters have made it a popular base for modification. Most
subsequent work adjusts the search behaviour: nonlinear or adaptive convergence
factors, chaotic initialisation, opposition-based learning, Gaussian and
L\'evy perturbations, crossover and elite mutation operators, and alternative
leadership selection rules \cite{...}. A second family combines GWO with local
search or with other metaheuristics in a hybrid scheme. A third family embeds
reinforcement learning, typically to select among search operators or to tune
parameters online \cite{...}. All of these modify \emph{how the wolves move}.

The present work is orthogonal to that axis. The movement equation of
\eqref{eq:move} is the published one; what is changed is which of the many
candidates the algorithm is allowed to verify with the true objective.

\subsection{Surrogate-assisted GWO for expensive optimization}

When evaluations are expensive, the established remedy is to replace the
objective with a cheap approximation and to call the true function only for a
few promising candidates. Surrogate-assisted GWO has been studied for
high-dimensional expensive black-box problems \cite{sagwo2020}, and
active-learning Kriging variants have followed. The approach is effective but
carries structural costs: training a Gaussian process or RBF model is
$O(N_s^3)$ in the number of samples, the model must be re-trained as the archive
grows, and approximation quality degrades in high dimensions. The candidates
chosen by an infill criterion are only as good as the model that ranks them.

RASE-GWO deliberately avoids this trade. It builds no model and performs no
matrix inversion; the screening rule uses only distances that are already
available from the movement step. The price is that the screening is purely
geometric and carries no estimate of fitness, which is acceptable precisely
because the quota is spent on candidates that are close to positions whose
fitness is known.

\subsection{Sparse and spatially screened evaluation}

A more recent line of work attacks the evaluation budget directly. Small-scale
learning PSO updates only one or two representative individuals per generation;
spatially screened PSO and sparse adaptive-quota CSO evaluate a fixed sparse
subset selected by geometric criteria and use the saved evaluations to fund
additional generations \cite{ssl-pso,ss-rlpso,saq-cso}. The motivation is the
same local-smoothness argument used here: once the swarm has contracted,
neighbouring individuals carry redundant information. Their empirical finding
that the evaluation quota should be biased strongly toward the promising region
motivates the action set used in Section~\ref{sec:q}.

Those frameworks were designed for algorithms whose individuals carry memory.
Transferring the same mechanism to GWO requires the verification anchor
introduced in Section~III: without it, the elite triangle can be occupied by
positions that were never evaluated. Furthermore, their exploration criterion
-- distance from the population centroid -- selects candidates that a
leader-driven algorithm cannot exploit, because a GWO candidate only influences
the search after entering the elite triangle. Section~\ref{sec:ablation}
quantifies both differences.

\subsection{Reinforcement learning for evaluation-budget management}

Reinforcement learning is widely used in evolutionary computation for parameter
control and operator scheduling \cite{...}. Using it to decide \emph{where the
evaluation budget goes} is less common, and in the sparse-evaluation
frameworks cited above the learned quantity is the schedule of search modes
rather than the allocation of evaluations among candidates. RASE-GWO learns the
split of the evaluation quota between consensus refinement and minority-leader
probing, which is a decision about the evaluations themselves. Because the
reward is bounded and scale invariant, a single Q-table is reused across
problems whose objective magnitudes differ by many orders of magnitude.
